\section*{Conclusiones del estado del arte}

Se observa que los NCA han evolucionado desde sus primeras aplicaciones en morfogénesis, donde se exploró su capacidad de generar estructuras complejas partiendo de una única célula activa, hacia modelos que buscan una mayor versatilidad generativa, capaces de producir dinámicas visuales, texturas, comportamientos adaptativos e incluso manipular estrategias, muchos de ellos inspirados en la definición de Mordvintsev. Además, se han desarrollado arquitecturas híbridas que no generan un NCA individual, sino que modelan un espacio de NCA, así como configuraciones generativas y bioinspiradas.

En conjunto, estos trabajos evidencian que los NCA no solo constituyen un marco computacional de interés teórico, sino también una herramienta con aplicaciones en inteligencia artificial, ciencias naturales e ingeniería. Esta diversidad de enfoques abre la puerta a explorar nuevas adaptaciones y propuestas, especialmente en áreas de aprendizaje profundo, autoorganización y simulación de procesos dinámicos.
